\section{Related Work}
\label{sec:related_work}

The interpretability of attention has been debated extensively in NLP. Jain and Wallace~\cite{jain2019attention} argue that standard attention weights often have weak correlation with gradient-based importance measures and that substantially different attention distributions can lead to similar predictions. Serrano and Smith~\cite{serrano2019attention} similarly show that large attention weights do not always correspond to large causal effects on model outputs. In response, Wiegreffe and Pinter~\cite{wiegreffe2019attention} argue that whether attention counts as an explanation depends on the definition of explanation and on how the evaluation is framed. Our work does not aim to settle this broader debate; instead, it tests a concrete and measurable version of the question in a simple Transformer classification setting.

We compare attention against two established families of attribution methods. Gradient-based methods estimate importance by differentiating the model score with respect to the input, and input $\times$ gradient is one of the best-known simple variants~\cite{ancona2018towards}. Perturbation-based methods estimate importance by observing how the prediction changes when part of the input is removed or altered. Representation erasure and leave-one-out style analyses have been used as a direct way to measure contribution by deletion~\cite{li2017erasure}. These two families provide useful comparison points because they capture different notions of importance: local sensitivity around the input versus explicit predictive change under removal.

Our setting is also related to analyses of attention patterns in pretrained Transformers. Clark et al.~\cite{clark2019bert} show that BERT attention often exhibits strong structural patterns, including frequent focus on delimiter tokens and position-based behaviors. Sun and Marasovi\'c~\cite{sun2021effective} further argue that raw attention can differ from more output-relevant effective attention. These observations motivate our decision to examine a simple last-layer attention signal empirically rather than assume that visual salience is equivalent to functional importance.
